\subsection{Results}

The normalized \ETPB and \EZDC distributions change systematically with $x_p$.
Relative to a low-$x_p$ selection, high-$x_p$ dijet events exhibit an \ETPB
distribution shifted toward lower energies. This is a direct confirmation of
the event-activity bias previously observed in centrality-dependent jet
measurements. The \EZDC distribution shows a qualitatively similar trend: at
high $x_p$, a larger fraction of events contains fewer forward neutrons. The
effect in the ZDC is, however, substantially weaker than that in the FCal.

The mean values \AZDC and \EATPB both decrease with increasing $x_p$, most
visibly for $x_p\gtrsim 0.02$, where earlier ATLAS measurements found evidence
for small proton-size configurations. Across this range, \AZDC decreases by up
to about $5\%$, corresponding on average to two fewer beam-energy neutrons,
whereas \EATPB decreases by up to $40\%$. When the two means are normalized to
their maximum values, the FCal estimator is found to be approximately six times
more sensitive to hard-process kinematics than the ZDC estimator.

The relation between \ETPB and \EZDC is also $x_p$ dependent. For ZDC energies
below approximately $150~\mathrm{TeV}$, the average FCal transverse energy is
consistent with a linear dependence on the ZDC energy in low-, middle-, and
high-$x_p$ selections. The slope decreases progressively from low to high
$x_p$, and the Pearson correlation coefficient changes by approximately $3\%$
across the measured range. At higher ZDC energies, deviations from the linear
trend are consistent with larger fluctuations in nuclear breakup for events
with the greatest event activity.

These results show that both common geometry estimators retain sensitivity to
the proton configuration associated with a hard scattering. The ZDC observable
is more robust than Pb-going FCal transverse energy, but is not independent of
the hard-scattering kinematics. This observation provides experimental input
for simultaneous descriptions of event activity, nuclear breakup, and color
fluctuations in $p$+Pb collisions.
